
\chapter{Battery Dynamics and Dispatch Problem}
\label{ch:battery-dynamics}

This chapter formalises the state-space model of a grid-scale battery energy
storage system (BESS), derives the dispatch optimisation programme, and
shows how observation uncertainty enlarges the feasible set in ways that
classical deterministic formulations cannot capture.

% ----------------------------------------------------------------
\section{Battery State Variables}

The BESS plant is characterised by the following state and parameter
variables at each hourly time step~$t$:

\begin{table}[ht]
\centering
\caption{Battery state and parameter variables.}
\label{tab:battery-vars}
\begin{tabular}{cll}
\toprule
\textbf{Symbol} & \textbf{Name} & \textbf{Units / Range} \\
\midrule
$E_t$          & State of charge (energy)      & MWh, $[0, E_{\max}]$ \\
$P_t^{+}$     & Charge power                  & MW,  $[0, P_{\max}^{+}]$ \\
$P_t^{-}$     & Discharge power               & MW,  $[0, P_{\max}^{-}]$ \\
$\eta^{+}$    & Charge efficiency              & dimensionless, $(0,1]$ \\
$\eta^{-}$    & Discharge efficiency           & dimensionless, $(0,1]$ \\
$E_{\max}$    & Nameplate energy capacity      & MWh \\
$P_{\max}$    & Rated power (symmetric)        & MW \\
$R_{\max}$    & Maximum ramp rate              & MW\,h$^{-1}$ \\
$\mathrm{SOC}_t$ & Normalised SOC, $E_t / E_{\max}$ & $[0, 1]$ \\
\bottomrule
\end{tabular}
\end{table}

For the ORIUS reference plant the nameplate parameters are
$E_{\max} = 10\,000$\,MWh, $P_{\max} = 200$\,MW, $\eta^{+} = \eta^{-} = 0.95$,
and $R_{\max} = P_{\max}$ (no additional ramp constraint beyond rated power).
Soft SOC margins are enforced at $[\mathrm{SOC}_{\min}, \mathrm{SOC}_{\max}]
= [0.05, 0.95]$; hard BMS trips occur at $[0.0, 1.0]$.

% ----------------------------------------------------------------
\section{SOC Dynamics}
\label{sec:soc-dynamics}

The energy balance over one time step ($\Delta t = 1$\,h) is:
\begin{equation}
\label{eq:soc-dynamics-battery}
  E_{t+1} = E_t + \eta^{+}\,P_t^{+}\,\Delta t
                 - \frac{P_t^{-}}{\eta^{-}}\,\Delta t.
\end{equation}
In normalised form:
\begin{equation}
\label{eq:soc-normalised}
  \mathrm{SOC}_{t+1} = \mathrm{SOC}_t
    + \frac{\eta^{+}\,P_t^{+} - P_t^{-}/\eta^{-}}{E_{\max}}\,\Delta t.
\end{equation}
This is a deterministic, linear transition---the stochasticity enters through
the \emph{observation} of $E_t$ (which may be stale or corrupted) and through
the forecast of external signals that drive the optimiser's choice of
$(P_t^{+}, P_t^{-})$.

% ----------------------------------------------------------------
\section{Power, Ramp-Rate, and SOC Constraints}
\label{sec:constraints}

The feasible action set at time~$t$ is defined by the following constraints:

\paragraph{Power bounds.}
\begin{equation}
\label{eq:power-bounds}
  0 \le P_t^{+} \le P_{\max}^{+}, \qquad
  0 \le P_t^{-} \le P_{\max}^{-}.
\end{equation}

\paragraph{Mutual exclusion.}
Simultaneous charge and discharge is prohibited:
\begin{equation}
\label{eq:mutual-exclusion}
  P_t^{+} \cdot P_t^{-} = 0.
\end{equation}
In the LP relaxation this is enforced via a binary indicator or by solving
separate charge and discharge sub-problems and selecting the dominant mode.

\paragraph{Ramp-rate limits.}
\begin{equation}
\label{eq:ramp-limits}
  |P_t^{+} - P_{t-1}^{+}| \le R_{\max}\,\Delta t, \qquad
  |P_t^{-} - P_{t-1}^{-}| \le R_{\max}\,\Delta t.
\end{equation}

\paragraph{SOC bounds.}
\begin{equation}
\label{eq:soc-bounds}
  E_{\min} \le E_{t+1} \le E_{\max},
\end{equation}
where $E_{\min} = \mathrm{SOC}_{\min} \cdot E_{\max}$ and similarly for the
upper bound.  Substituting the dynamics~\eqref{eq:soc-dynamics-battery}:
\begin{equation}
\label{eq:soc-bounds-expanded}
  E_{\min} - E_t \;\le\;
    \eta^{+}\,P_t^{+}\,\Delta t - \frac{P_t^{-}}{\eta^{-}}\,\Delta t
  \;\le\; E_{\max} - E_t.
\end{equation}

% ----------------------------------------------------------------
\section{Candidate, Safe, and Executed Actions}

Three action vectors are distinguished at each control step:

\begin{description}
  \item[$a_t^\star = (P_t^{+\star}, P_t^{-\star})$:]
    The \emph{candidate action} produced by the dispatch optimiser, maximising
    expected economic value subject to the forecast distribution and the
    \emph{observed} SOC~$\tilde{E}_t$.

  \item[$a_t^{\mathrm{safe}} = (P_t^{+,\mathrm{safe}}, P_t^{-,\mathrm{safe}})$:]
    The \emph{safe action} produced by the DC3S shield
    (Chapter~\ref{ch:dc3s-adapter}).  This is the nearest feasible point to
    $a_t^\star$ within the tightened safe-action set~$\mathcal{A}_t$.

  \item[$a_t^{\mathrm{exec}}$:]
    The \emph{executed action} sent to the BMS.  Under nominal operation
    $a_t^{\mathrm{exec}} = a_t^{\mathrm{safe}}$; under BMS override
    $a_t^{\mathrm{exec}}$ may be further clipped by hardware limits.
\end{description}

The \emph{intervention rate} (IR) measures the fraction of steps where
$a_t^{\mathrm{safe}} \neq a_t^\star$, quantifying the economic cost of
safety enforcement.

% ----------------------------------------------------------------
\section{Dispatch Objective}
\label{sec:dispatch-objective}

The deterministic dispatch LP over a rolling horizon~$H$ is:
\begin{equation}
\label{eq:dispatch-lp}
\begin{aligned}
  \max_{P_{t:t+H}^{+},\, P_{t:t+H}^{-}} \quad &
    \sum_{\tau=t}^{t+H-1}
    \bigl[\,\hat{\pi}_\tau\,(P_\tau^{-} - P_\tau^{+})
           - c_{\mathrm{deg}}\,(P_\tau^{+} + P_\tau^{-})\,\bigr]\,\Delta t \\
  \text{s.t.} \quad &
    E_{\tau+1} = E_\tau + \eta^{+}\,P_\tau^{+}\,\Delta t
                       - P_\tau^{-}/\eta^{-}\,\Delta t,
      \quad \forall\,\tau, \\
  & E_{\min} \le E_{\tau+1} \le E_{\max},
      \quad \forall\,\tau, \\
  & 0 \le P_\tau^{+} \le P_{\max}^{+},\;
    0 \le P_\tau^{-} \le P_{\max}^{-},
      \quad \forall\,\tau, \\
  & |P_\tau^{+} - P_{\tau-1}^{+}| \le R_{\max}\,\Delta t,
      \quad \forall\,\tau, \\
  & |P_\tau^{-} - P_{\tau-1}^{-}| \le R_{\max}\,\Delta t,
      \quad \forall\,\tau, \\
  & E_t = \tilde{E}_t \quad (\text{initial condition from observation}).
\end{aligned}
\end{equation}
Here $\hat{\pi}_\tau$ is the forecast price, $c_{\mathrm{deg}}$ is a
degradation penalty per MWh throughput, and the initial SOC~$\tilde{E}_t$ is
the \emph{observed} (possibly degraded) state.

The objective trades revenue (buy low, sell high) against degradation cost.
Only the first action $(P_t^{+\star}, P_t^{-\star})$ is committed; the
remainder is discarded and re-solved at $t+1$ (receding-horizon control).

% ----------------------------------------------------------------
\section{Uncertainty and Feasible Set Under Degraded Observation}

When the observation $\tilde{E}_t$ is degraded, the initial-condition
constraint $E_t = \tilde{E}_t$ in~\eqref{eq:dispatch-lp} becomes uncertain.
Let $\hat{E}_t$ denote the controller's best estimate and
$\delta_t = E_t - \hat{E}_t$ the estimation error.  The true SOC after action
is:
\begin{equation}
\label{eq:soc-uncertain}
  E_{t+1}^{\mathrm{true}} = (\hat{E}_t + \delta_t)
    + \eta^{+}\,P_t^{+}\,\Delta t - \frac{P_t^{-}}{\eta^{-}}\,\Delta t.
\end{equation}
For the SOC bounds to hold with probability at least $1 - \alpha$, the
constraint must be tightened by the conformal prediction half-width~$q_t$:
\begin{equation}
\label{eq:tightened-bounds}
  E_{\min} + q_t \;\le\; E_{t+1}^{\mathrm{est}} \;\le\; E_{\max} - q_t,
\end{equation}
where $q_t$ is the reliability-inflated interval half-width from RAC-Cert
(Chapter~\ref{ch:dc3s-adapter}).  This tightening shrinks the feasible set
and reduces the economic value of the dispatch, but guarantees safety
against observation error.

% ----------------------------------------------------------------
\section{Calibration Gap and Uncertainty Mismatch}

Even with conformal intervals, a \emph{calibration gap} arises when marginal
coverage differs from conditional coverage across reliability groups.  The
locked CQR group-coverage results
(\texttt{reports/publication/cqr\_group\_coverage.csv}) reveal that wind
intervals achieve only 39--49\% coverage in the low and high reliability
groups despite 90\% nominal marginal coverage.  Solar intervals are better
calibrated (82--90\%) but still exhibit under-coverage at the high-variability
group.

This calibration gap means that the tightened bounds
in~\eqref{eq:tightened-bounds} are \emph{insufficiently tight} during periods
of poor observation quality, motivating the adaptive inflation law in DC3S
(Section~\ref{sec:dispatch-objective} cross-references
Chapter~\ref{ch:dc3s-adapter}).

% ----------------------------------------------------------------
\section{Baseline Policies}

The ORIUS evaluation includes eight dispatch controllers spanning four
design philosophies:

\begin{enumerate}
  \item \textbf{Deterministic LP} (\texttt{deterministic\_lp}):
    solves~\eqref{eq:dispatch-lp} on point forecasts with no interval
    awareness.
  \item \textbf{Robust fixed-interval} (\texttt{robust\_fixed\_interval}):
    worst-case robust LP using fixed (non-adaptive) prediction intervals.
  \item \textbf{CVaR interval} (\texttt{cvar\_interval}):
    CVaR-optimised scenario dispatch over quantile forecasts.
  \item \textbf{ACI conformal} (\texttt{aci\_conformal}):
    adaptive conformal inference baseline with online width adjustment.
  \item \textbf{Scenario robust} (\texttt{scenario\_robust}):
    scenario-based robust optimisation over sampled forecast paths.
  \item \textbf{Scenario MPC} (\texttt{scenario\_mpc}):
    receding-horizon model predictive control with scenario tree.
  \item \textbf{DC3S wrapped} (\texttt{dc3s\_wrapped}):
    full DC3S shield with linear reliability inflation.
  \item \textbf{DC3S FTIT} (\texttt{dc3s\_ftit}):
    DC3S with fault-tolerant invariant tube extension.
\end{enumerate}

Results across these controllers are reported in
Chapter~\ref{ch:dc3s-adapter} and the benchmark track
(Chapter~\ref{ch:cpsbench}).

% ----------------------------------------------------------------
\section{Scope Boundary}

The dispatch model in this dissertation is a \emph{single-asset, energy-only}
formulation:

\begin{itemize}
  \item \textbf{Included:} energy arbitrage, SOC management, degradation
    penalty, safety constraints.
  \item \textbf{Excluded:} ancillary services (frequency regulation, spinning
    reserve), multi-asset portfolio dispatch, network constraints, and
    market-clearing co-optimisation.
\end{itemize}

The single-asset scope is sufficient to demonstrate the observation-quality
feedback mechanism; extension to fleet-level dispatch is discussed in
Chapter~\ref{ch:orius-bench-battery}.

% ----------------------------------------------------------------
\section{Chapter Summary}

This chapter formalised the battery state-space model, the SOC dynamics
equation~\eqref{eq:soc-dynamics-battery}, the power/ramp/SOC constraint set, and the
deterministic dispatch LP~\eqref{eq:dispatch-lp}.  Three action types---candidate,
safe, and executed---were distinguished to precisely locate the safety
intervention point.  The effect of degraded observation on the feasible set
was derived, showing that the conformal half-width must tighten the SOC bounds
and that calibration gaps in wind and solar intervals motivate adaptive
inflation.  Eight baseline controllers spanning deterministic, robust,
conformal, and shielded design philosophies were enumerated.
